\documentclass{beamer}
\usetheme{Antibes}
\usepackage{xcolor, colortbl}
\usepackage{algorithm}
\usepackage[noend]{algpseudocode}
\usepackage{textcomp}
\usepackage{listings}
\usepackage{hyperref}
\usepackage{alltt}
\usepackage{tikz}
\usepackage{framed}
\usepackage{marvosym}
\usepackage{wasysym}
\usepackage{marvosym}
\usepackage{crayola}
\usepackage{mathpartir}
\usepackage{tabularx}
\usepackage[belowskip=-15pt,aboveskip=0pt]{caption}
\usepackage[skins]{tcolorbox}
\usepackage{multicol}
\usetikzlibrary{positioning,shapes,arrows, backgrounds, fit, shadows}
\usetikzlibrary{decorations.markings}
%\usepackage{wasysym}
%\usepackage{marvosym}
\setbeamertemplate{footline}[frame number]
%\usecolortheme{fly}
\usefonttheme{serif}

\title[Sujit]{Programming Language Implementation -- Introduction}
\author{Sujit Kumar Chakrabarti}
\institute{IIITB}
\date{IIIT-B}

%\begin{document}
%\setbeamertemplate{navigation symbols}{}%remove navigation symbols
%
%\maketitle

\definecolor{lightblue}{rgb}{0.8,0.93,1.0} % color values Red, Green, Blue
\definecolor{darkblue}{rgb}{0.4,0.3,1.0} % color values Red, Green, Blue
\definecolor{Blue}{rgb}{0,0,1.0} % color values Red, Green, Blue
\definecolor{darkgreen}{rgb}{0,0.7,0.2} % color values Red, Green, Blue
\definecolor{Red}{rgb}{1,0,0} % color values Red, Green, Blue
\definecolor{Pink}{rgb}{0.7,0,0.2}
\definecolor{links}{HTML}{2A1B81}
\definecolor{mydarkgreen}{HTML}{126215}
\newcommand{\highlight}[1]{{\color{Red}(#1)}}

\newcommand{\myheader}[1]{
	{\color{darkblue}
		\begin{Large}
			\begin{center}
				{#1}
			\end{center}
		\end{Large}
	}
}
\newcommand{\myminorheader}[1]{
	{\color{BrickRed}
		\begin{Large}
			{\fontfamily{\sfdefault}\selectfont\textbf{#1}}
		\end{Large}
	}
}

%\tikzstyle{input} = [coordinate]
%\tikzstyle{output} = [coordinate]


\tikzstyle{bb}=[%
      rectangle, draw=black, thick, fill=OliveGreen!30, drop shadow, align=center,
      text ragged, minimum height=2em, minimum width=2em, inner sep=6pt
]

\tikzstyle{inv}=[%
      rectangle, draw=none,  align=center,
      text ragged, minimum height=2em, minimum width=2em, inner sep=6pt
]

\tikzstyle{db}=[%
      ellipse, draw=black, thick, fill=pink, drop shadow, align=center,
      text ragged, minimum height=2em, inner sep=6pt
]

\tikzstyle{jn}=[%
      ellipse, draw=black, thick, fill=black
]

\tikzstyle{io}=[%
      trapezium, trapezium left angle=60, trapezium right angle=120, draw=black, thick, fill=brown, drop shadow,
      text ragged, minimum height=2em, minimum width=2em, inner sep=6pt, align=center
]

\tikzstyle{glio}=[%
      trapezium, trapezium left angle=60, trapezium right angle=120, draw=red, line width = 1mm, fill=brown, drop shadow,
      text ragged, minimum height=2em, minimum width=2em, inner sep=6pt
]
\tikzstyle{gl}=[%
      rectangle, draw=red, line width = 1mm, fill=lightblue, drop shadow,
      text ragged, minimum height=2em, minimum width=2em, inner sep=6pt
]

\tikzstyle{en}=[%
      rectangle, draw=black, thick, fill=none,
      text ragged, minimum height=2em, minimum width=2em, inner sep=6pt
]

\tikzstyle{sh}=[%
      rectangle, draw=gray, thick, fill=none, color = gray,
      text ragged, minimum height=2em, minimum width=2em, inner sep=6pt
]


\lstdefinestyle{javacode}{
	language = Java,
	basicstyle = \ttfamily\scriptsize,
	stringstyle = \ttfamily,
	keywordstyle=\color{Blue}\bfseries,
	identifierstyle=\color{Pink},
	commentstyle=\color{darkgreen},
	frame=single,
	frameround=tttt,
%	numbers=left
	showstringspaces=false
}

\lstdefinestyle{camlcode}{
	language = Caml,
	basicstyle = \scriptsize\ttfamily,
	stringstyle = \color{red}\ttfamily,
	keywordstyle=\color{Blue}\bfseries,
	identifierstyle=\ttfamily,
	frame=single,
	frameround=tttt,
	numbers=none,
	showstringspaces=false,
	escapeinside={(*@}{@*)}
}

\lstdefinestyle{outputcode}{
	language = bash,
	backgroundcolor = \color{black},
	basicstyle = \tiny\ttfamily\color{white},
	stringstyle = \color{red}\ttfamily,
	keywordstyle=\color{white}\bfseries,
	identifierstyle=\ttfamily,
	frameround=tttt,
	numbers=none,
	showstringspaces=false,
	escapeinside={(*@}{@*)}
}

\newtcolorbox{myframe}[2][]{%
  enhanced,colback=white,colframe=black,coltitle=black,
  sharp corners,boxrule=0.4pt,
  fonttitle=\itshape,
  attach boxed title to top left={yshift=-0.3\baselineskip-0.4pt,xshift=2mm},
  boxed title style={tile,size=minimal,left=0.5mm,right=0.5mm,
    colback=white,before upper=\strut},
  title=#2,#1
}


\begin{document}
\maketitle

\definecolor{lightblue}{rgb}{0.8,0.93,1.0} % color values Red, Green, Blue
\definecolor{darkblue}{rgb}{0.4,0.3,1.0} % color values Red, Green, Blue
\definecolor{Blue}{rgb}{0,0,1.0} % color values Red, Green, Blue
\definecolor{darkgreen}{rgb}{0,0.7,0.2} % color values Red, Green, Blue
\definecolor{Red}{rgb}{1,0,0} % color values Red, Green, Blue
\definecolor{Pink}{rgb}{0.7,0,0.2}
\definecolor{links}{HTML}{2A1B81}
\definecolor{mydarkgreen}{HTML}{126215}

\mathchardef\mhyphen="2D

\newsavebox{\prt}
\savebox{\prt}{
  \begin{tikzpicture}
    \draw[line width=0.5mm] (0, 0) -- (.2, -.2) -- (0, -.4);
  \end{tikzpicture}
}



\newsavebox{\sta}
\savebox{\sta}{
\resizebox{0.4\textwidth}{!}{%
  \begin{tikzpicture}
	\node (1) {$=$};
    \node (2) [below left = of 1] {$<id, 1>$};
    \node (3) [below right = of 1] {$+$};
    \node (4) [below left = of 3] {$<id, 2>$};
    \node (5) [below right = of 3] {$*$};
    \node (6) [below left = of 5] {$<id, 3>$};
    \node (7) [below right = of 5] {$60$};
    \path (1) edge  (2)
          (1) edge  (3)
          (3) edge  (4)
          (3) edge  (5)
          (5) edge  (6)
          (5) edge  (7)
    ;
    
  \end{tikzpicture}
}}

\newsavebox{\stb}
\savebox{\stb}{
\resizebox{0.4\textwidth}{!}{%
  \begin{tikzpicture}
	\node (1) {$=$};
    \node (2) [below left = of 1] {$<id, 1>$};
    \node (3) [below right = of 1] {$+$};
    \node (4) [below left = of 3] {$<id, 2>$};
    \node (5) [below right = of 3] {$*$};
    \node (6) [below left = of 5] {$<id, 3>$};
    \node (7) [below right = of 5] {\textbf{inttofloat}};
    \node (8) [below       = of 7] {$60$};

    \path (1) edge  (2)
          (1) edge  (3)
          (3) edge  (4)
          (3) edge  (5)
          (5) edge  (6)
          (5) edge  (7)
          (7) edge  (8)
    ;
    
  \end{tikzpicture}
}}

\tikzstyle{init} = [draw=black, thick, fill=black, circle]
\tikzstyle{tr} = [rectangle,draw=black, very thin, dashed,align=left,font=\scriptsize, text ragged, minimum height=1em, minimum width=2em, inner sep=1pt, outer sep=3pt]
\tikzstyle{trl} = [->, thick, color=Purple, text=black]
\tikzset{obj/.style = {rectangle split, rounded corners,
		rectangle split parts=2, thick,draw=black, top
		color=Dandelion!10,bottom color=Dandelion!20, align=center}}

\tikzset{cobj/.style = {rectangle split, rounded corners,
		rectangle split parts=3, thick,draw=black, top
		color=white,bottom color=Dandelion!10, align=center}}

\tikzset{sn/.style = {rectangle, draw=gray, fill=blue!10, thick, align=center, rounded corners}}
\tikzset{sdn/.style = {ellipse, draw=gray, fill=red!10, thick, align=center}}
\tikzset{st/.style = {-latex, draw=brown, thick}}

\newcommand\Mstrut[1]{\rule{0pt}{#1cm}}


\tikzstyle{ss} = [draw=gray, thick, fill=white, rounded corners, rectangle, inner sep=15pt]

% frame begin %%%%%%%%%%%%%%%%%%%%%%%%
\begin{frame}{Compilers}
{Overview}

\begin{center}

\begin{tikzpicture}[auto,
    ->,
  %  shorten >=2pt,
    >=stealth,
  %  node distance=1cm,
]
    \node[inv] (0) []             {Source \\ Code};
    \node[bb]  (1) [right = of 0] {Compiler \\ (Translator)};
    \node[inv] (2) [right = of 1] {Output \\ Code};

    \path (0) edge (1)
          (1) edge (2)
    ;

  \end{tikzpicture}
\end{center}

\begin{itemize}
\item Translates input program to output program
\item Input program (source code ) in high-level language (e.g. C, Java, Python etc.)
\item Output program (target code) in low-level language (e.g. machine code, assembly, byte-code etc.)
\item Source code $\equiv$ Target code
\end{itemize}
\end{frame}
% frame end %%%%%%%%%%%%%%%%%%%%%%%%

% frame begin %%%%%%%%%%%%%%%%%%%%%%%%
\begin{frame}{Compiler Structure}

\begin{center}
\resizebox{\textwidth}{!}{%
\begin{tikzpicture}[auto,
    ->,
  %  shorten >=2pt,
    >=stealth,
  %  node distance=1cm,
]
    \node[inv] (0) []             {Source \\ Code ($L_1$)};
    \node[bb]  (1) [right = of 0] {Front \\ End};
    \node[bb]  (2) [right = of 1] {Back \\ End};
    \node[inv] (3) [right = of 2] {Output \\ Code ($L_2$)};
\begin{small}
    \path (0) edge (1)
          (1) edge node[align=center] {IR}             (2)
          (2) edge              (3)
    ;

\end{small}
  \end{tikzpicture}
}

\end{center}
\begin{itemize}
\item Two functions:
	\begin{enumerate}
	\item Read source code and analyse for errors
	\item Translate into the target code
	\end{enumerate}
\item Two parts:
	\begin{enumerate}
	\item Front end
	\item Back end
	\end{enumerate}
\pause
\item \textbf{Intermediate representation (IR)}
	\begin{itemize}
	\item  A data structure output from front end and input to the backend
	\item Several phases; several IRs
	\end{itemize}
\end{itemize}

\end{frame}
% frame end %%%%%%%%%%%%%%%%%%%%%%%%


% frame begin %%%%%%%%%%%%%%%%%%%%%%%%
\begin{frame}{Compiler Structure}
{Front End}
\begin{center}
\resizebox{\textwidth}{!}{%
\begin{tikzpicture}[auto,
    ->,
  %  shorten >=2pt,
    >=stealth,
  %  node distance=1cm,
]
    \node[inv] (inp) []             {Source \\ Code};
    \node[bb]  (synan) [right = of inp] {Syntax \\ Analyser};
    \node[bb]  (seman) [right = of synan] {Semantic \\ Analyser};
    \node[bb]  (opt) [right = of seman] {Optimiser};
    \node[bb]  (cg) [right = of opt] {Code \\ Generator};
    \node[inv] (outp) [right = of cg] {Output \\ Code};

\node[dashed, thick, draw=Gray, fit=(synan) (seman), minimum height=3cm] (fe) {};
\node[inv, above=0.1cm of fe.south](lfe) {\textbf{\textsf{\color{Brown}Front End}}};

\node[dashed, thick, draw=Gray, fill=Gray, fit=(opt) (cg), minimum height=3cm] (be) {};
\node[inv, above=0.1cm of be.south](lbe) {\textbf{\textsf{\color{Brown}Back End}}};

\begin{small}
\draw[->] (inp) -- (synan);
\draw[->] (synan) -- node[above]{IR1}(seman);
\draw[->] (seman) -- node[above]{IR2}(opt);
%\draw[->] (opt) -- node[above]{IR3}(cg);
\draw[->] (cg) -- (outp);
\end{small}
  \end{tikzpicture}
}
\end{center}


\end{frame}
% frame end %%%%%%%%%%%%%%%%%%%%%%%%

% frame begin %%%%%%%%%%%%%%%%%%%%%%%%
\begin{frame}{Compiler Structure}
{Syntax Analysis}

\begin{itemize}
\item \textbf{\color{BrickRed}Syntax errors:}
	\begin{enumerate}
	\item Example 1: \lstinline[style=javacode]|for i in range(10)|
	\item Example 2: \lstinline[style=javacode]|void f(x) \{ ... \}|
	\end{enumerate}
\item When program doesn't follow grammatical constructs of the input PL
\item \textbf{Well formed programs:} Programs with \emph{no} syntax errors
\item \textbf{Ill formed programs:} Programs \emph{with} syntax errors
\end{itemize}
\end{frame}
% frame end %%%%%%%%%%%%%%%%%%%%%%%%

% frame begin %%%%%%%%%%%%%%%%%%%%%%%%
\begin{frame}{Compiler Structure}
{Syntax Analysis}

\begin{center}
\resizebox{!}{0.5\textheight}{%
\begin{tikzpicture}[auto,
    ->,
  %  shorten >=2pt,
    >=stealth,
  %  node distance=1cm,
]
    \node[inv] (1)                {\lstinline@pos = init + r * 60@ \\ \tiny{$<$id, 1$>$ $<$=$>$ $<$id, 2$>$ $<$+$>$ $<$id, 3$>$ $<$*$>$ $<$60$>$}};
    \node[bb]  (2) [below = of 1] {Syntax Analysis};
    \node[inv] (3) [below = of 2] {\usebox{\sta}};

    \path (1) edge  (2)
          (2) edge  (3)
    ;
  \end{tikzpicture}
}

\end{center}

\pause
\begin{itemize}
\item Successful syntax analysis $\longrightarrow$ Abstract syntax tree (AST)
\item More details later
\end{itemize}
\end{frame}
% frame end %%%%%%%%%%%%%%%%%%%%%%%%

% frame begin %%%%%%%%%%%%%%%%%%%%%%%%
\begin{frame}{Compiler Structure}
{Lexical Analyser}
\begin{center}
\resizebox{\textwidth}{!}{%
\begin{tikzpicture}[auto,
    ->,
  %  shorten >=2pt,
    >=stealth,
  %  node distance=1cm,
]
    \node[inv] (inp) []             {Source \\ Code};
    \node[bb]  (lexan) [right = of inp] {Lexical \\ Analyser};
    \node[bb]  (synan) [right = of lexan] {Syntax \\ Analyser};
    \node[bb]  (seman) [right = of synan] {Semantic \\ Analyser};
    \node[bb]  (opt) [right = of seman] {Optimiser};
    \node[bb]  (cg) [right = of opt] {Code \\ Generator};
    \node[inv] (outp) [right = of cg] {Output \\ Code};

\node[dashed, thick, draw=Gray, fit=(lexan) (seman), minimum height=3cm] (fe) {};
\node[inv, above=0.1cm of fe.south](lfe) {\textbf{\textsf{\color{Brown}Front End}}};

\node[dashed, thick, draw=Gray, fill=Gray, fit=(opt) (cg), minimum height=3cm] (be) {};
\node[inv, above=0.1cm of be.south](lbe) {\textbf{\textsf{\color{Brown}Back End}}};

\begin{small}
\draw[->] (inp) -- (lexan);
\draw[->] (lexan) -- node[above]{tokens}(synan);
\draw[->] (synan) -- node[above]{IR1}(seman);
\draw[->] (seman) -- node[above]{IR2}(opt);
%\draw[->] (opt) -- node[above]{IR3}(cg);
\draw[->] (cg) -- (outp);
\end{small}
  \end{tikzpicture}
}
\end{center}


\end{frame}
% frame end %%%%%%%%%%%%%%%%%%%%%%%%

% frame begin %%%%%%%%%%%%%%%%%%%%%%%%
\begin{frame}{Compiler Structure}
{Lexical Analysis}

\begin{center}
\resizebox{\textwidth}{!}{%
\begin{tikzpicture}[auto,
    ->,
  %  shorten >=2pt,
    >=stealth,
  %  node distance=1cm,
]
    \node[inv] (1)                {\lstinline@pos = init + r * 60@};
    \node[bb]  (2) [below = of 1] {Lexical Analysis};
    \node[inv] (3) [below = of 2] {$<$id, 1$>$ $<$=$>$ $<$id, 2$>$ $<$+$>$ $<$id, 3$>$ $<$*$>$ $<$60$>$};
\begin{small}

    \path (1) edge  (2)
          (2) edge  (3)
    ;

\end{small}
  \end{tikzpicture}
}

\end{center}

\end{frame}
% frame end %%%%%%%%%%%%%%%%%%%%%%%%

% frame begin %%%%%%%%%%%%%%%%%%%%%%%%
\begin{frame}[fragile]{Compiler Structure}
{Semantic Analysis -- Functions}

\begin{enumerate}
\item Variable binding:
\begin{lstlisting}[style=camlcode]
let x = 20 in
let sum1 = 2 * x in
let x = 100 and z = 5 in
  sum1 - z + x
\end{lstlisting}
\item Type checking:
\begin{lstlisting}[style=camlcode]
void add(int x, int y) { return x + y; }
add(1, 2);
\end{lstlisting}
\item Type checking:
\begin{lstlisting}[style=camlcode]
# let add x y = x + y;;
val add : int -> int -> int = <fun>
\end{lstlisting}
\end{enumerate}
\end{frame}
% frame end %%%%%%%%%%%%%%%%%%%%%%%%

% frame begin %%%%%%%%%%%%%%%%%%%%%%%%
\begin{frame}[fragile]{Compiler Structure}
{Semantic Analysis -- Semantic errors}

\begin{itemize}
\item \textbf{\color{BrickRed}Semantic errors:} Error associated with the meaning of the program
\item Natural language example:
\begin{quote}
Table goes to school.
\end{quote}
\pause
\item \textbf{Ill-typed programs:} Well formed programs that don't make sense.
\item Examples:
	\begin{enumerate}
	\item Example 1: \lstinline[style=javacode]|String s = 10;|
	\item Example 2: \lstinline[style=javacode]|void f() \{ return x; \}|
	\end{enumerate}
\end{itemize}
\end{frame}
% frame end %%%%%%%%%%%%%%%%%%%%%%%%

% frame begin %%%%%%%%%%%%%%%%%%%%%%%%
\begin{frame}{Compiler Structure}
{Semantic Analysis}


\begin{center}
\resizebox{!}{0.7\textheight}{%
\begin{tikzpicture}[auto,
    ->,
  %  shorten >=2pt,
    >=stealth,
  %  node distance=1cm,
]
    \node[inv] (1)                {\usebox{\sta}};
    \node[bb]  (2) [below = of 1] {Semantic Analyser};
    \node[inv] (3) [below = of 2] {\usebox{\stb}};

    \path (1) edge  (2)
          (2) edge  (3)
    ;
  \end{tikzpicture}
}

\pause
\begin{tikzpicture}
\node[rectangle, draw=Red, fill=Red, minimum height=0.25cm](){};
\end{tikzpicture}

\end{center}

\end{frame}
% frame end %%%%%%%%%%%%%%%%%%%%%%%%


% frame begin %%%%%%%%%%%%%%%%%%%%%%%%
\begin{frame}{Compiler Structure}
{Next Module}
\myheader{Compiler Backend}

\end{frame}
% frame end %%%%%%%%%%%%%%%%%%%%%%%%

% frame begin %%%%%%%%%%%%%%%%%%%%%%%%
\begin{frame}{Compiler Structure}
{Back End}
\begin{center}
\resizebox{\textwidth}{!}{%
\begin{tikzpicture}[auto,
    ->,
  %  shorten >=2pt,
    >=stealth,
  %  node distance=1cm,
]
    \node[inv] (inp) []             {Source \\ Code};
    \node[bb]  (synan) [right = of inp] {Syntax \\ Analyser};
    \node[bb]  (seman) [right = of synan] {Semantic \\ Analyser};
    \node[bb]  (opt) [right = of seman] {Optimiser};
    \node[bb]  (cg) [right = of opt] {Code \\ Generator};
    \node[inv] (outp) [right = of cg] {Output \\ Code};

\node[dashed, thick, draw=Gray, fill=Gray, fit=(synan) (seman), minimum height=3cm] (fe) {};
\node[inv, above=0.1cm of fe.south](lfe) {\textbf{\textsf{\color{Brown}Front End}}};

\node[dashed, thick, draw=Gray, fit=(opt) (cg), minimum height=3cm] (be) {};
\node[inv, above=0.1cm of be.south](lbe) {\textbf{\textsf{\color{Brown}Back End}}};

\begin{small}
\draw[->] (inp) -- (synan);
%\draw[->] (synan) -- node[above]{IR1}(seman);
\draw[->] (seman) -- node[above]{IR2}(opt);
\draw[->] (opt) -- node[above]{IR3}(cg);
\draw[->] (cg) -- (outp);
\end{small}
  \end{tikzpicture}
}
\end{center}


\end{frame}
% frame end %%%%%%%%%%%%%%%%%%%%%%%%

% frame begin %%%%%%%%%%%%%%%%%%%%%%%%
\begin{frame}{Compiler Structure}
{Code Optimisation}

\begin{center}
\resizebox{!}{0.8\textheight}{%
\begin{tikzpicture}[auto,
    ->,
  %  shorten >=2pt,
    >=stealth,
]
    \node[inv, align=left] (1){
    		\lstinline@t1 = inttofloat(60)@ \\ 
    		\lstinline@t2 = id3 * t1@ \\
    		\lstinline@t3 = id2 + t2@ \\ 
    		\lstinline@id1 = t3@
	};
    \node[bb]  (2) [below = of 1] {Code Optimiser};
    \node[inv] (3) [below = of 2] {
    		\lstinline@t1 = id3 * 60.0@ \\
    		\lstinline@id1 = id2 + t1@
	};
\begin{small}

    \path (1) edge  (2)
          (2) edge  (3)
    ;

\end{small}
  \end{tikzpicture}
}

\end{center}

\end{frame}
% frame end %%%%%%%%%%%%%%%%%%%%%%%%

% frame begin %%%%%%%%%%%%%%%%%%%%%%%%
\begin{frame}[fragile]{Compiler Structure}
{Code Optimisation}

\myheader{Example -- Loop optimisation}


\begin{center}
\begin{tabular}{c @{\hspace{0.5cm}$\Rightarrow$\hspace{0.5cm}} c} 
\begin{minipage}{0.45\textwidth}
\begin{lstlisting}[language=java, basicstyle=\scriptsize\ttfamily]
for(int i = 0; i < N; i++) {
  t = a + b;
  sum += t + i;
}
\end{lstlisting}
\end{minipage}
&
\begin{minipage}{0.45\textwidth}
\begin{lstlisting}[language=java, basicstyle=\scriptsize\ttfamily]
t = a + b;
for(int i = 0; i < N; i++) {
  sum += t + i;
}
\end{lstlisting}
\end{minipage}
\end{tabular}
\end{center}

\end{frame}
% frame end %%%%%%%%%%%%%%%%%%%%%%%%

% frame begin %%%%%%%%%%%%%%%%%%%%%%%%
\begin{frame}{Compiler Structure}
{Code Optimisation}

\begin{itemize}
\item Numerous algorithms
\item Local, global, interprocedural
\item Control flow analysis, data flow analysis
\item Machine independent and machine dependent
\item ...
\pause
\item Largest and most complex module of modern compilers
\item Active area of research
\end{itemize}

\end{frame}
% frame end %%%%%%%%%%%%%%%%%%%%%%%%

% frame begin %%%%%%%%%%%%%%%%%%%%%%%%
\begin{frame}{Compiler Structure}
{Code Generation}

\begin{center}
\resizebox{!}{0.8\textheight}{%
\begin{tikzpicture}[auto,
    ->,
  %  shorten >=2pt,
    >=stealth,
]
    \node[inv, align=left] (1){
    		\lstinline@t1 = id3 * 60.0@ \\
    		\lstinline@id1 = id2 + t1@
	};
    \node[bb]  (2) [below = of 1] {Code Generator};
    \node[inv, align=left] (3) [below = of 2] {
    		\lstinline@LDF R2, id3@ \\
    		\lstinline@MULF R2, R2, \#60.0@ \\
    		\lstinline@LDF R1, id2@ \\
    		\lstinline@ADDF R1, R1, R2@ \\
    		\lstinline@STF id1, R1@ \\
	};
\begin{small}

    \path (1) edge  (2)
          (2) edge  (3)
    ;

\end{small}
  \end{tikzpicture}
}
\end{center}

\end{frame}
% frame end %%%%%%%%%%%%%%%%%%%%%%%%

% frame begin %%%%%%%%%%%%%%%%%%%%%%%%
\begin{frame}{Compiler Structure}
{Code Generation}

\myheader{Runtime Organisation}
\begin{center}
\resizebox{0.6\textwidth}{!}{%
\begin{tikzpicture}[auto,
    ->,
  %  shorten >=2pt,
    >=stealth,
]
\begin{small}
\node[inv](lstack){\textsc{\textbf{\underline{Stack}}}};
\node[rectangle, draw=Black, thick, minimum width=1cm, below=0.1cm of lstack] (a) {a};
\node[rectangle, draw=Gray, dashed, thick, fit=(lstack)(a)] (stack) {};


\node[inv, right=1cm of lstack](lheap){\textsc{\textbf{\underline{Heap}}}};
\node[rectangle, draw=Black, thick, minimum width=1cm, below=0.1cm of lheap] (o) {o};
\node[rectangle, draw=Gray, dashed, thick, fit=(lheap)(o), minimum height=2cm] (heap) {};

\node[inv, right=1.5cm of lheap](lstatic){\textsc{\textbf{\underline{Static}}}};

\node[inv, below left=0.1cm of lstatic, xshift=0.5cm](lA){\textsc{\textbf{\underline{A}}}};
\node[rectangle, draw=Black, thick, minimum width=1cm, below=0.1cm of lA] (Am1) {$m_1$};
\node[rectangle, draw=Black, thick, minimum width=1cm, below=0.1cm of Am1] (Am2) {$m_2$};
\node[rectangle, draw=Brown, thick, fit=(lA) (Am2)] (A) {};

\node[inv, below right=0.1cm of lstatic, xshift=-0.5cm](lB){\textsc{\textbf{\underline{B}}}};
\node[rectangle, draw=Black, thick, minimum width=1cm, below=0.1cm of lB] (Bm1) {$m_1$};
\node[rectangle, draw=Black, thick, minimum width=1cm, below=0.1cm of Bm1] (Bm2) {$m_2$};
\node[rectangle, draw=Black, thick, minimum width=1cm, below=0.1cm of Bm2] (Bm3) {$m_3$};

\node[rectangle, draw=Brown, thick, fit=(lB) (Bm3)] (B) {};

\node[rectangle, draw=Gray, dashed, thick, fit=(lstatic) (A) (B)] (static) {};

\node[inv, right=1.5cm of lstatic](lcode){\textsc{\textbf{\underline{Code}}}};
\node[rectangle, draw=Black, thick, minimum width=1cm, below=0.1cm of lcode] (procAm1) {$A::m_1$};
\node[rectangle, draw=Black, thick, minimum width=1cm, below=0.1cm of procAm1] (procAm2) {$A::m_2$};
\node[rectangle, draw=Black, thick, minimum width=1cm, below=0.1cm of procAm2] (procBm1) {$B::m_1$};
\node[rectangle, draw=Black, thick, minimum width=1cm, below=0.1cm of procBm1] (procBm2) {$B::m_2$};
\node[rectangle, draw=Black, thick, minimum width=1cm, below=0.1cm of procBm2] (procBm3) {$B::m_3$};
\node[rectangle, draw=Gray, dashed, thick, fit=(lcode) (procBm3)] (code) {};

\draw[->, dotted] (a) -- (o);
\draw[->, dotted, out=-90, in=-90] (o.south) to (B.south);
\draw[->, dotted, in=0] (Am1) -- (procAm1);
\draw[->, dotted, in=0] (Am2) -- (procAm2);
\draw[->, dotted, in=0] (Bm1) -- (procBm1);
\draw[->, dotted, in=0] (Bm2) -- (procBm2);
\draw[->, dotted, in=0] (Bm3) -- (procBm3);

\end{small}
  \end{tikzpicture}
}
\end{center}

\begin{scriptsize}
\begin{itemize}
\item Placement of various parts of the target code in the appropriate parts of the virtual memory
\item Not everything generated statically
\item Certain things have to be created by the target code at runtime.
\end{itemize}
\end{scriptsize}
\end{frame}
% frame end %%%%%%%%%%%%%%%%%%%%%%%%

% frame begin %%%%%%%%%%%%%%%%%%%%%%%%
\begin{frame}{Compiler Structure}

\myheader{Register allocation}

\begin{center}
\includegraphics[width=0.8\textwidth]{images/reg-alloc.png}
\end{center}

\end{frame}
% frame end %%%%%%%%%%%%%%%%%%%%%%%%

% frame begin %%%%%%%%%%%%%%%%%%%%%%%%
\begin{frame}{Compiler Structure}

\begin{center}
\resizebox{\textwidth}{!}{%
\begin{tikzpicture}[auto,
    ->,
  %  shorten >=2pt,
    >=stealth,
  %  node distance=1cm,
]
    \node[inv] (0) []             {};
    \node[bb]  (1) [right = of 0, xshift=1cm] {Lexical \\ Analysis};
    \node[bb]  (2) [right = of 1, xshift=1cm] {Syntax \\ Analysis};
    \node[bb]  (3) [right = of 2, xshift=1cm] {Semantic \\ Analysis};
    \node[bb]  (4) [right = of 3, xshift=1cm] {Optimisation};
    \node[bb]  (5) [right = of 4, xshift=1cm] {Code \\ Generation};
    \node[inv] (6) [right = of 5, xshift=1cm] {};
\begin{small}
    \path (0) edge node[align=center] {Source \\ Code}     (1)
          (1) edge node[align=center] {Tokens}             (2)
          (2) edge node[align=center] {Abstract \\ Syntax} (3)
          (3) edge node[align=center] {Annotated \\ IR}    (4)
          (4) edge node[align=center] {Transformed \\ IR}  (5)
          (5) edge node[align=center] {Generated \\ Code}  (6)
    ;

\end{small}
  \end{tikzpicture}
}

\end{center}

\end{frame}
% frame end %%%%%%%%%%%%%%%%%%%%%%%%

% frame begin %%%%%%%%%%%%%%%%%%%%%%%%
\begin{frame}{Summary -- PLDI}
{Applications}
\begin{itemize}
	\item Hardware synthesis
	\item Database query interpreter
	\item IDEs
	\item Debuggers
	\item Bug finders -- Lint, Valgrind
\end{itemize}

\end{frame}
% frame end %%%%%%%%%%%%%%%%%%%%%%%%

% frame begin %%%%%%%%%%%%%%%%%%%%%%%%
\begin{frame}{Summary -- PLDI}
{Why study PLDI}
\begin{itemize}
	\item Vibrant area of study -- new PLs everyday!
	\item Sophisticated theory + Practical system building
	\item Deeper knowledge of data-structures, algorithms, software design
\end{itemize}

\begin{center}
\pause
\begin{tikzpicture}
\node[rectangle, draw=Red, fill=Red, minimum height=0.25cm](){};
\end{tikzpicture}
\end{center}
\end{frame}
% frame end %%%%%%%%%%%%%%%%%%%%%%%%

\end{document}
